\chapter{Conclusões}
\label{cap:conclusao}

\todo{adicionar conclusao relacionada ao eixo extensao, com os resultados e concluindo que isso justifica a necessidade de haver um ids baseado em redes neurais para detecçao}

Todas as atividades planejadas no cronograma do plano de trabalho foram executadas, e algumas foram estendidas além do escopo original: a atualização do levantamento bibliográfico, o estudo do protocolo DNS (incluindo captura prática de tráfego com Wireshark), a investigação de \textit{datasets} de tráfego DNS (CIC-Bell-DNS2021 e BCCC-CIC-Bell-DNS-2024), o estudo de redes neurais e de arquiteturas de IDS, a criação e avaliação de um novo modelo de classificação (SVM, KNN e MLP, em cenários binário e multiclasse, com e sem a \textit{feature} \texttt{src\_port}, e com balanceamento de classes), a especialização e comparação do modelo para domínios brasileiros (.br), a avaliação de arquiteturas híbridas de \textit{deep learning} (\textit{ensemble} MLP+CNN+LSTM, entre outras), a implementação de um pipeline de captura e classificação de domínios em tempo real, e a análise e disseminação dos resultados por meio de publicação científica e disponibilização de código aberto.

O principal resultado deste trabalho foi a criação de um sistema capaz de identificar domínios maliciosos a partir do tráfego DNS, especialmente para o contexto brasileiro. Além disso, contribuiu-se com a identificação de um viés de aprendizado (a \textit{feature} \texttt{src\_port}) presente no dataset BCCC-CIC-Bell-DNS-2024, cuja remoção e substituição por \textit{features} genuinamente comportamentais permitiu construir um modelo MLP com desempenho elevado e mais robusto para a tarefa de detecção de tráfego DNS malicioso, além de evidenciar, na prática, a importância da explicabilidade de modelos de aprendizado de máquina aplicados a sistemas de detecção de intrusão. Complementarmente, o projeto avançou além do previsto no plano de trabalho ao validar a viabilidade de um modelo especializado para domínios \texttt{.br}, ao explorar arquiteturas híbridas de \textit{deep learning} e ao implementar um pipeline funcional de classificação de domínios a partir de captura de tráfego DNS real.

Como trabalhos futuros, ficam propostos: (i) a coleta ou construção de um \textit{dataset} genuinamente nativo de domínios e tráfego DNS brasileiros, e não apenas uma filtragem de amostras \texttt{.br} dentro de um dataset internacional, de forma a permitir uma validação mais robusta da especialização do modelo para padrões de \textit{typosquatting} do contexto nacional; (ii) a exploração de técnicas de \textit{Neural Architecture Search}~\citep{ren2021,zoph2016} e de seleção automática de modelos por \textit{meta-learning}~\citep{sylligardos2023} para otimizar de forma mais sistemática as arquiteturas híbridas já testadas; (iii) a complementação do pipeline de integração já implementado com a geração automática de regras de bloqueio em servidores DNS ou \textit{firewalls} a partir da classificação obtida, completando o terceiro item do procedimento de mitigação proposto no artigo~\citep{hirata2026dns}; e (iv) a avaliação do modelo e do pipeline em um ambiente de rede real, em produção, com monitoramento contínuo de desempenho e reforço periódico do treinamento (\textit{retraining}) à medida que novos padrões de ataque surgirem.